\section{Vanilla Tradition}

\paragraph{Living Dangerously}


There are two directions for future developments of Tradition, based on how seriously one takes Guenon's view of the modern world. Specifically, he regarded it as an ``illusion'', i.e., totally false, and it would dissipate as soon as the illusion evaporated. That means he opposed its very foundations, including the ideas of progress, evolution, democracy, equality, etc., that is, every quality that the modern world holds dear. From his perspective, these ideas are not worth debating; only an ``intellectual conversion'' could destroy the illusion.

Guenon began publishing his books in the 1920s. Apart from Guido De Giorgio who knew him personally, it seems that Julius Evola was the first writer to have incorporated Guenon's views into a larger whole. That was still in the 20s, well before figures like Ananda Coomaraswamy, Fritjhof Schuon became involved. Schuon, in particular, has had a minor following who consider themselves the guardians of Tradition. They are the nice side of tradition, earned by writing favorably about world religions, even compiling a compendium of ``Traditional Wisdom''. They often aim for academic respectability. This ``vanilla'' tradition seems far from the radical nature of Guenon himself. Moreover, this strand tends to exclude Evola from their purview, mostly by neglect, yet resorting to hostility when necessary.

Let us not forget that Guenon regarded the diremption of politics from spiritual life as another deviation of the modern world. Hence, the ``separation of church and state'' or the French-style secular state are simply errors promulgated in the modern world. Now it is true that apart from his involvement with Action Française, which ended with its condemnation by Pio IX, Guenon avoided active participation in politics, for obvious reasons. Evola, on the other hand, again for the obvious reason, did engage in the political situation at his time. Of course, Evola made some mistakes, since action in the world differs from the contemplation of metaphysical principles.

This hardly means that he was anti-Tradition, as the ``nice'' wing would have it. Quite the contrary, he was following one of Guenon's core principles. Now who else, among the new breed of professors and professional ``traditionalists'', are so intransigently opposed to the values and mores of the modern world?

\paragraph{Path of Blame}


\textbf{Charles Upton}, in \emph{The System of the Antichrist}, gives a useful summary of the ideal of the People of Blame\footnote{\url{https://en.wikipedia.org/wiki/Malamatiyya}}. The true esoterist becomes indistinguishable from exoteric believers. Ibn al-Arabi pointed out that the spiritual works which are obligatory for ordinary believers are even more so obligatory to the Sufi esoterist. The principle of non-duality is the metaphysical justification for this: God is beyond the duality of esoteric/exoteric, so one is not ``higher'' than the other.

This is the same conclusion that Valentin Tomberg reached: anthroposophy cannot be an alternative to the Traditional church. Authority comes through Tradition, not from one's particular spiritual experiences. The latter must be subservient to the former.

This became crystal clear to me some time ago, while I was participating in a weekly discussion group on \textbf{Nisargadatta}'s\footnote{\url{http://sri-nisargadatta-maharaj.blogspot.com}} discourse collected in \emph{I Am That}. Since he was regarded as an enlightened being (I am not disputing that), we strove to penetrate his mystical utterances. Of course, we were all ``beyond'' any exoteric religions. One day, someone brought in a documentary film about him. I noticed right away that he had a little altar set up in his apartment. So, certainly, such a man does not neglect his exoteric prayers. Once the depth of that observation sunk in, I left the group shortly thereafter.

Now Nisargadatta seemed to be a humble man, running his tobacco shop, and not running expensive weekend seminars.

\paragraph{Sudden Illumination}


That is not the case, however, for most New Age gurus. Eckhart Tolle\footnote{\url{http://www.eckharttolle.com}}, for example, followed a rather formulaic path. At the point of despair, he came to a sudden inner transformation leading to the realization of his life purpose. His reading of Bo Yin Ra\footnote{\url{https://www.gornahoor.net/?p=4611}} was the trigger.

That worked pretty well for Mr Tolle. Unfortunately, it won't work for you. You will have to pay through the nose for books, videos, overpriced lectures and expensive retreats, wherein you will experience various ``modalities'' to dissolve the ego. As if the ego will spend all that money just to get dissolved! But it is not just about dissolving the ego, since a side effect will be the elimination of violence in the world. Any day now.


\flrightit{Posted on 2015-02-15 by Cologero}

\begin{center}* * *\end{center}

\begin{footnotesize}
\begin{sffamily}
\texttt{Jonathan Olavsson on 2015-02-16 at 00:59 said:}

Responding to the poll question:

I would greatly appreciate to read `Sintesi di Dottrina della Razza'. But not merely because it is of historical interest -- even because it is my understanding that Evola's spiritual approach to the question of race came very close to the actual truth. I will learn Italian, so the book will be read eventually even if an English translation isn't made available. But such a translation would doubtless be an enlightening resource for numerous Anglophones.

\hfill

\texttt{Cologero on 2015-02-16 at 02:47 said:}

I can't imagine what you might mean by ``numerous'', Mr Olavsson. A few years ago, I offered the translation to a mailing list. All I asked for in return was an intelligent discussion on the list. I couldn't get 50 people to sign up ... and that's out of 5 to 10 thousand visitors to this site per month. And there is more material even beyond Evola which I am gathering.

\hfill

\texttt{Jonathan Olavsson on 2015-02-16 at 03:29 said:}

To those worth counting, it should be of great interest. How many those are, you are probably better qualified at estimating.

Which material are you gathering?

\hfill

\texttt{rhondda9 on 2015-02-16 at 10:26 said:}

I am very interested in this, but I am afraid that I am very slow and have to contemplate what I have read for a period of time. I am more of a reader than an active participant in a discussion. Usually I have to put what I have read in my own words and then hope they are not foolish.

\hfill

\texttt{Tom Blanchard on 2015-02-16 at 10:48 said:}

To be fair, Gornahoor has readers today that it did not have years ago. Judging by the quality of the comments on recent articles, I think it may be a good time.

Working on the Maserati, Cologero, as I can't seem to find a copy of ``La Razza di Roma'' for sale anywhere. What interior/exterior colors / options were you thinking of?

\hfill

\texttt{Boreas on 2015-02-16 at 16:52 said:}

I certainly would like to be able to read `Sintesi... ' in English, and certainly not as a historical curiosity.

\hfill

\texttt{Nexist418 on 2015-02-17 at 22:20 said:}

I've been wanting to read Sintesi for some time now.

\hfill

\texttt{Cologero on 2015-02-17 at 23:27 said:}

You can find all currently available texts here: \url{http://www.gornahoor.net/?tag=sintesi-di-dottrina-della-razza}.

I will be encouraged by thoughtful comments.

\hfill

\texttt{ozzy belcim on 2015-02-18 at 13:19 said:}

Cologero, first of all, the work you are doing on this blog and the material you cover are truly commendable in a cyberspace crippled with a dearth of thought. However, you should be the first to know that just like the material, the audience should be elite by definition, therefore few. I can't imagine a time past or future when multitudes can interest themselves in the issues at stake here. Therefore, don't let the numbers upset you. It is bound to be that way. As a researcher into the spread of ideas in the public space, I can tell you that your indicated 10,000 monthly visits does not really correspond to more than 200- 300 unique monthly visitors. I don't want to dwell much on the fact that a substantial portion of these are just passers-by. According to website statistics such as alexa, average visitor to your site spends 2,5 minutes here, which is hardly sufficient to read, comprehend and ponder on the ideas treated here. I am a subscriber to nielsen bookscan globally and I can tell you the same dynamic at play all around the world. You would be surprised how few serious books are sold in US and elsewhere. All Guenon titles published and sold in US over 15 years suggest these books don't reach more than a few thousand people. I can say the same thing -- although 10-20 times those of Guenon -- for even writers such as Kant and Heidegger in the profane domain in spite of the number of university departments and students studying them. The advent of e-readers, electronic copies and ipads etc. has no impact on the historical sales figures of serious books. It seems your average kindle reader riding the subway is not a market for such. This is after all what Guenon has in mind when he talks about the intellectual elite although today's populace is too brainwashed and programmed against the idea of any elite based on some spurious democratic sensibility. However, intellectual elitism seems to be ingrained in human society and is indisputably demonstrable even with quantitative analysis, which the modern world is so enamored of.

\hfill

\texttt{Cologero on 2015-02-26 at 00:26 said:}

Good points, Ozzy. Gornahoor does not subscribe to Alexa, nor to Google's enhanced demographics. They won't benefit me, just advertisers. You are correct, a lot of ``visitors'' are just bots and spammers trying to insert comments. Guenon has something more in mind beyond mere intellectual capability. He calls for an ``intellectual conversion'' which is a step beyond reading serious books. This bears repeating.

As an aside, I was at a tech event this week with several Indian Hindus in attendance. I asked each one of them if they had ever read the Bhagavad Gita. I got just one positive response out of ten. I'm not quite sure of what they thought of me as I tried to explain the BG to them as best I could! It seems that they are satisfied with a statue of Ganesh in their cars.

\hfill

\texttt{Cologero on 2015-02-26 at 00:28 said:}

Miser, \textit{The Turd in the Punch Bowl}\footnote{\url{http://www.gornahoor.net/?p=7354}} was for your benefit.

\hfill

\texttt{Cologero on 2015-02-26 at 12:35 said:}

Nero Ribelle, Mr Blanchard, with a contrasting interior.

Whatever happened to ``free speech'' when certain groups can suppress books they deem unflattering?

\end{sffamily}
\end{footnotesize}

